\documentclass{TesiMatteoBianchi}

\titolo{Progettazione e sviluppo della gestione dei task in sistemi multi-agente basati sul protocollo A2A e LangGraph}
\autore{Matteo Bianchi}
\relatore{Prof. Maurizio Patrignani}
\correlatore{Ing. Fabio Barillari \\ Dott. Federico Chinotto}


\annoAccademico{2025/2026}

\dedica{Questa è la dedica}

\usepackage[plainpages=false,hypertexnames=false]{hyperref}
\usepackage[all]{hypcap}
\usepackage{xurl} 
\usepackage{amsthm}
\usepackage{amsmath}
\usepackage{tikz}
\usetikzlibrary{arrows.meta, positioning, decorations.pathreplacing}
\usepackage{algorithm}
\usepackage{algpseudocode}
\usepackage{placeins}
\algdef{SE}[TRY]{Try}{EndTry}{\textbf{try}}{\textbf{end try}}
\algdef{C}[TRY]{TRY}{Catch}[1]{ \textbf{catch} #1}
\algdef{C}[TRY]{TRY}{Finally}{\textbf{finally}}

\let\underscoreoriginale\_
\renewcommand{\_}{\underscoreoriginale\allowbreak}
\setlength{\emergencystretch}{3em}


\begin{document}
\frontmatter
\generaFrontespizio
\generaDedica
\ringraziamenti{ringraziamenti}
\introduzione{introduzione}
\generaIndice

\mainmatter
\capitolo{Dagli LLM agli agenti}{capitoli/capitolo1_llm-e-agenti}
\capitolo{Il protocollo Agent2Agent}{capitoli/capitolo2_protocollo-a2a}
\capitolo{Requisiti del Sistema}{capitoli/capitolo3_requisiti-sistema}
\capitolo{Progettazione}{capitoli/capitolo4_progettazione}

\capitolo{Validazione e Testing}{capitoli/capitolo5_validazione-testing}



\backmatter
\conclusioni{conclusioni}

\bibliography{bibliografia}

\end{document}